\subsection{Sequences, Cauchy Criteria, and Completeness}

% **Role:** rigorous language of real and complex analysis.
%
% Limits of real sequences, subsequences, Cauchy sequences, monotone convergence, limsup/liminf, Bolzano--Weierstrass.

Sequences are the first testing ground for limits, order, and completeness. The main tools are subsequences, Cauchy criteria, monotone convergence, limsup/liminf, and Bolzano--Weierstrass.

\begin{definition} \label{def:real-complex-sequence-convergence} \hlt{Convergent Real and Complex Sequences}\\
A real or complex sequence $(s_n)$ \hlt{converges} to $s$ if, $\forall \varepsilon>0$, $\exists N\in\N$ s.t. $n\geq N$ implies $\abs{s_n-s}<\varepsilon$.\\
Notation is $s_n\to s$ or $\lim\limits_{n\to\infty}s_n=s$. If no such $s$ exists, the sequence \hlt{diverges}.
\end{definition}

\begin{figure}[H]
  \centering
  \begin{tikzpicture}[
  >=Stealth,
  x=0.9cm,
  y=0.9cm,
  every node/.style={font=\small},
  axis/.style={->, gray!70, line width=0.55pt},
  early/.style={circle, fill=gray!65, inner sep=1.35pt},
  late/.style={circle, fill=blue!70!black, inner sep=1.45pt},
  limit/.style={circle, fill=orange!85!black, inner sep=1.55pt},
  interval/.style={blue!65!black, line width=0.8pt}
]
  \draw[axis] (0,0) -- (8.0,0);

  \coordinate (L) at (5.8,0);
  \coordinate (Left) at (4.75,0);
  \coordinate (Right) at (6.85,0);

  \draw[interval] (Left) -- (Right);
  \node[gray!70] at (Left) {$($};
  \node[gray!70] at (Right) {$)$};
  \node[below] at (Left) {$s-\varepsilon$};
  \node[below] at (Right) {$s+\varepsilon$};
  \node[limit] at (L) {};
  \node[below] at (L) {$s$};

  \foreach \x/\lab in {0.9/s_1,1.75/s_2,2.60/s_3} {
    \node[early] at (\x,0) {};
    \node[below] at (\x,0) {\(\lab\)};
  }
  \node[gray!70] at (3.30,0) {$\cdots$};

  \foreach \x in {4.98,5.26,5.52,6.02,6.32,6.60} {
    \node[late] at (\x,0) {};
  }
\end{tikzpicture}

  \caption{For all $n\geq N$, the terms $s_n$ stay inside $(s-\varepsilon,s+\varepsilon)$.}
  \label{fig:sequence-epsilon-neighborhood}
\end{figure}

\begin{definition} \label{def:real-sequence-range-and-boundedness} \hlt{Range and Boundedness of a Real Sequence}\\
The \hlt{range} of $(s_n)$ is set $\{s_n:n\in\N\}$. Real sequence is \hlt{bounded} if its range is bounded above and below in $\R$.
\end{definition}

\begin{theorem} \label{thm:complex-sequence-limit-algebra} \hlt{Limit Algebra for Complex Sequences}\\
Suppose $(s_n)$ and $(t_n)$ are complex sequences with $s_n\to s$ and $t_n\to t$. Then, for every $c\in\C$,
\begin{enumerate}[label=(\roman*)]
\setlength{\itemsep}{0pt}
\item $s_n+t_n\to s+t$;
\item $cs_n\to cs$ and $c+s_n\to c+s$;
\item $s_nt_n\to st$;
\item if $s\neq0$ and $s_n\neq0$ for all $n$, then $1/s_n\to1/s$.
\end{enumerate}
\begin{proof}
\begin{enumerate}[label=(\roman*)]
\setlength{\itemsep}{0pt}
\item Given $\varepsilon>0$, choose $N$ such that $\abs{s_n-s}<\varepsilon/2$, $\abs{t_n-t}<\varepsilon/2$ for $n\geq N$.\\
Then $\abs{(s_n+t_n)-(s+t)}\leq\abs{s_n-s}+\abs{t_n-t}<\varepsilon$.
\item Scalar multiplication uses $\abs{cs_n-cs}=\abs{c}\abs{s_n-s}$; translation is (i) with constant sequence $c$.
\item Since $s_n\to s$, choose $N\geq2$ with $\abs{s_n-s}<1$ for $n\geq N$ and set $M\defeq\max\{\abs{s_1},\ldots,\abs{s_{N-1}},\abs{s}+1\}$.\\
Then $\abs{s_nt_n-st}\leq \abs{s_n}\abs{t_n-t}+\abs{t}\abs{s_n-s}\leq M\abs{t_n-t}+\abs{t}\abs{s_n-s}\to0$.
\item If $s\neq0$, then for all sufficiently large $n$, $\abs{s_n}>\abs{s}/2$. Hence $\abs{\frac1{s_n}-\frac1s}=\frac{\abs{s_n-s}}{\abs{s_n}\abs{s}}\leq \frac{2}{\abs{s}^2}\abs{s_n-s}\to0$.
\end{enumerate}
\hlt{Structural Commentary}: Boundedness controls products; nonzero limits control reciprocals.\\
\hlt{Key Insight}: Product and reciprocal rules are estimates after the easy sum rule.
\end{proof}
\end{theorem}

\begin{theorem} \label{thm:rk-sequence-coordinate-convergence} \hlt{Coordinatewise Convergence in $\R^k$}\\
Let $\mathbf{x}_n=(\alpha_{j,n})_{j=1}^k$ and $\mathbf{x}=(\alpha_j)_{j=1}^k$ in $\R^k$. Then $\mathbf{x}_n\to\mathbf{x}$ $\Longleftrightarrow$ $\alpha_{j,n}\to\alpha_j$ $(1\leq j\leq k)$.\\
Moreover, if $\mathbf{x}_n\to\mathbf{x}$, $\mathbf{y}_n\to\mathbf{y}$, and $\beta_n\to\beta$ in $\R$, then $\mathbf{x}_n+\mathbf{y}_n\to\mathbf{x}+\mathbf{y}$; $\mathbf{x}_n\cdot\mathbf{y}_n\to\mathbf{x}\cdot\mathbf{y}$; $\beta_n\mathbf{x}_n\to\beta\mathbf{x}$.
\begin{proof}
\begin{enumerate}[label=\roman*.]
\setlength{\itemsep}{0pt}
\item[$(\Rightarrow)$] Since $\abs{\alpha_{j,n}-\alpha_j}\leq\norm{\mathbf{x}_n-\mathbf{x}}_2$ for each $j$, each coordinate converges.
\item[$(\Leftarrow)$] Given $\varepsilon>0$, choose $N$ so that $\abs{\alpha_{j,n}-\alpha_j}<\varepsilon/\sqrt{k}$ for all $j$ and $n\geq N$. Then $\norm{\mathbf{x}_n-\mathbf{x}}_2<\varepsilon$.
\end{enumerate}
The algebraic conclusions follow coordinatewise from Theorem \ref{thm:complex-sequence-limit-algebra}.\\
\hlt{Structural Commentary}: Euclidean convergence is exactly simultaneous coordinate convergence.\\
\hlt{Key Insight}: The norm packages finitely many coordinate estimates into one estimate.
\end{proof}
\end{theorem}

\begin{remark} \label{rem:coordinate-metric-convergence-equivalence} \hlt{Coordinate Metrics Give the Same Limits}\\
For $k\in\N$ and $x,y\in\R^k$, set $d_1(x,y)=\sum\limits_{j=1}^k\abs{x_j-y_j}$, $d_2(x,y)=\big(\sum\limits_{j=1}^k\abs{x_j-y_j}^2\big)^{1/2}$,  $d_\infty(x,y)=\max\limits_{1\leq j\leq k}\abs{x_j-y_j}$.\\
Since $d_\infty\leq d_2\leq d_1\leq k d_\infty$, these have same convergent sequences, equivalently same coordinatewise criterion.
\end{remark}

\begin{definition} \label{def:subsequence-and-subsequential-limit} \hlt{Subsequences and Subsequential Limits}\\
Let $(p_n)$ be a sequence. If $n_1<n_2<n_3<\cdots$ is a strictly increasing sequence of positive integers, then $(p_{n_k})$ is a \hlt{subsequence} of $(p_n)$. If $p_{n_k}\to p$, then $p$ is a \hlt{subsequential limit} of $(p_n)$.
\end{definition}

\begin{remark} \label{rem:subsequence-selection-map} \hlt{Subsequences as Selection Maps}\\
Equivalently, a subsequence is the composite of the original sequence, viewed as a function on $\N$ as in Definition \ref{def:sequence-as-enumeration}, with a strictly increasing map $\sigma:\N\to\N$. In this notation the subsequence is $k\mapsto p_{\sigma(k)}$.\\
\hlt{Key Insight}: A subsequence keeps the original order but is allowed to skip terms.
\end{remark}

\begin{remark} \label{rem:subsequences-of-convergent-sequences} \hlt{Subsequences of Convergent Sequences}\\
If $p_n\to p$, then every subsequence of $(p_n)$ also converges to $p$. Conversely, if every subsequence of $(p_n)$ converges to $p$, then $(p_n)$ itself converges to $p$, since $(p_n)$ is one of its subsequences.
\end{remark}

\begin{definition} \label{def:monotone-sequence} \hlt{Monotone Real Sequence}\\
A real sequence $(s_n)$ is \hlt{monotonically increasing} if $s_n\leq s_{n+1}$ for all $n$, and \hlt{monotonically decreasing} if $s_n\geq s_{n+1}$ for all $n$. A \hlt{monotonic} sequence is either increasing or decreasing.
\end{definition}

\begin{definition} \label{def:dominant-index} \hlt{Dominant Index}\\
For a real sequence $(s_n)$, an index $n$ is \hlt{dominant} if $s_n\geq s_m$ for every $m>n$.
\end{definition}

\begin{theorem} \label{thm:monotone-subsequence-theorem} \hlt{Monotone Subsequence Theorem}\\
Every real sequence has a monotone subsequence.
\begin{proof}
If the dominant indices $n_1<n_2<\cdots$ are infinite, then $s_{n_1}\geq s_{n_2}\geq\cdots$, giving a decreasing subsequence.\\
Otherwise choose $N$ beyond all dominant indices and set $n_1>N$.\\
Given $n_k>N$, nondominance gives $n_{k+1}>n_k$ with $s_{n_{k+1}}>s_{n_k}$; hence $(s_{n_k})$ is increasing.\\
\hlt{Structural Commentary}: Dominant indices give decreasing case; nondominance recursively gives increasing case.\\
\hlt{Key Insight}: The proof organizes arbitrary real behavior into one ordered thread.
\end{proof}
\end{theorem}

\begin{figure}[H]
  \centering
  \begin{tikzpicture}[
  >=Stealth,
  x=0.9cm,
  y=0.9cm,
  every node/.style={font=\small},
  axis/.style={<->, gray!70, line width=0.5pt},
  point/.style={circle, fill=black, inner sep=1.35pt},
  snpoint/.style={circle, fill=blue!70!black, inner sep=1.45pt},
  supoint/.style={circle, fill=green!50!black, inner sep=1.45pt}
]
  \draw[axis] (0.5,0) -- (7.55,0);

  \coordinate (L) at (4.45,0);
  \coordinate (S) at (5.65,0);
  \coordinate (R) at (6.85,0);
  \coordinate (SN) at (4.74,0);

  \node[gray!70] at (L) {$($};
  \node[gray!70] at (R) {$)$};

  \draw[green!50!black, dashed, line width=0.45pt] (S) -- ++(0,0.62);
  \node[below] at (L) {$s-\varepsilon$};
  \node[below] at (R) {$s+\varepsilon$};
  \node[green!50!black, anchor=west] at ($(S)+(0,0.78)$) {$s=\sup\{s_n:n\in\N\}$};

  \foreach \x/\lab in {1.75/s_1,2.50/s_2,3.25/s_3} {
    \node[point] at (\x,0) {};
    \node[above] at (\x,0.12) {$\lab$};
  }
  \node[above] at (3.90,0.12) {$\cdots$};

  \foreach \x in {5.08,5.30,5.44,5.53,5.59,5.63} {
    \node[point] at (\x,0) {};
  }
  \node[snpoint] at (SN) {};
  \node[supoint] at (S) {};

  \draw[blue!70!black, dashed, line width=0.45pt] (SN) -- ++(0,0.62);
  \node[blue!70!black] at ($(SN)+(0,0.78)$) {$s_N$};
\end{tikzpicture}

  \caption{For an increasing bounded sequence, $s=\sup\{s_n\}$ is the limit.}
  \label{fig:monotone-convergence-supremum}
\end{figure}

\begin{theorem} \label{thm:monotone-convergence} \hlt{Monotone Convergence Theorem for Sequences}\\
A monotonic real sequence converges if and only if it is bounded.
\begin{proof}
\begin{enumerate}[label=\roman*.]
\setlength{\itemsep}{0pt}
\item[$(\Rightarrow)$] If $s_n\to s$, choose $N$ with $\abs{s_n-s}<1$ for $n\geq N$. The tail lies in $(s-1,s+1)$ and the initial set $\{s_1,\ldots,s_{N-1}\}$ is finite, so $(s_n)$ is bounded.
\item[$(\Leftarrow)$] Let $E\defeq\{s_n:n\in\N\}$. If $(s_n)$ is increasing, set $s\defeq\sup E$. Given $\varepsilon>0$, choose $N$ with $s-\varepsilon<s_N\leq s$; then $s-\varepsilon<s_n\leq s$ for $n\geq N$, so $s_n\to s$. If $(s_n)$ is decreasing, set $s\defeq\inf E$. Given $\varepsilon>0$, choose $N$ with $s\leq s_N<s+\varepsilon$; then $s\leq s_n\leq s_N<s+\varepsilon$ for $n\geq N$, so $s_n\to s$.
\end{enumerate}
\hlt{Structural Commentary}: Monotonicity traps bounded tails at $\sup E$ or $\inf E$.\\
\hlt{Key Insight}: The least upper bound is forced to be the limit.
\end{proof}
\end{theorem}

\begin{theorem} \label{thm:bolzano-weierstrass-real-sequence} \hlt{Bolzano--Weierstrass for Real Sequences}\\
Every bounded real sequence has a convergent subsequence.
\begin{proof}
By Theorem \ref{thm:monotone-subsequence-theorem}, choose a monotone subsequence. It is bounded because the original sequence is bounded, so Theorem \ref{thm:monotone-convergence} gives convergence.\\
\hlt{Structural Commentary}: Boundedness prevents escape; monotone subsequences force one ordered thread.\\
\hlt{Key Insight}: Monotone subsequence plus monotone convergence is the real-line compactness mechanism.
\end{proof}
\end{theorem}

\begin{definition} \label{def:real-cauchy-sequence} \hlt{Cauchy Real Sequence}\\
A real sequence $(s_n)$ is \hlt{Cauchy} if, for every $\varepsilon>0$, there is $N\in\N$ such that $m,n\geq N$ implies $\abs{s_n-s_m}<\varepsilon$.
\end{definition}

\begin{lemma} \label{lem:real-cauchy-sequence-with-convergent-subsequence} \hlt{Cauchy Plus One Convergent Subsequence}\\
If a real Cauchy sequence has a subsequence converging to $s\in\R$, then the whole sequence converges to $s$.
\begin{proof}
Given $\varepsilon>0$, choose $N_0$ so that $m,n\geq N_0$ implies $\abs{s_m-s_n}<\varepsilon/2$. Choose $j$ with $n_j\geq N_0$ and $\abs{s_{n_j}-s}<\varepsilon/2$. Then for every $n\geq N_0$, $\abs{s_n-s}\leq\abs{s_n-s_{n_j}}+\abs{s_{n_j}-s}<\varepsilon$.\\
\hlt{Structural Commentary}: Cauchy tails follow any convergent thread.\\
\hlt{Key Insight}: Mutual tail closeness turns one subsequential limit into the full limit.
\end{proof}
\end{lemma}

\begin{theorem} \label{thm:real-cauchy-convergence-criterion} \hlt{Real Cauchy Criterion}\\
A real sequence converges in $\R$ if and only if it is Cauchy.
\begin{proof}
If $s_n\to s$, choose $N$ such that $\abs{s_n-s}<\varepsilon/2$ for $n\geq N$; then $\abs{s_n-s_m}<\varepsilon$ for $m,n\geq N$. Conversely, suppose $(s_n)$ is Cauchy. Choose $N$ so that $m,n\geq N$ implies $\abs{s_m-s_n}<1$. Then the tail lies in $(s_N-1,s_N+1)$, and the finitely many initial terms are bounded, so $(s_n)$ is bounded. By Theorem \ref{thm:bolzano-weierstrass-real-sequence}, some subsequence converges; Lemma \ref{lem:real-cauchy-sequence-with-convergent-subsequence} makes the whole sequence converge.\\
\hlt{Structural Commentary}: Completeness of $\R$ means Cauchy behavior has a real endpoint.\\
\hlt{Key Insight}: Cauchy gives boundedness, Bolzano--Weierstrass gives a subsequential endpoint, and the Cauchy property pulls the whole sequence there.
\end{proof}
\end{theorem}
